\section{Results}
\label{sec:results}

\begin{figure*}[t]
\centering
\includegraphics[width=\textwidth]{figures/fdr_all.pdf}
\caption{Test-time FDR (mean$\pm$std over 100 splits) across all 32 configurations. Rows correspond to the two candidate models; columns to the four datasets. Each line represents a different judge (see legend). The dashed diagonal marks FDR~$= \alpha$.}
\label{fig:fdr_all}
\end{figure*}

We evaluate whether the framework delivers on its two central promises: (1)~the FDR guarantee holds empirically across judges, candidates, and datasets, and (2)~adaptive retrieval recovers coverage that would otherwise be lost to abstention. We use predictive entropy and LLM evaluator for the main results, and additionally analyze the sensitivity of our framework to the choice of admission function and retrieval depth (\S\ref{sec:sensitivity}).

\subsection{FDR control holds across all configurations}
\label{sec:fdr_results}

The most basic requirement of our framework is that the empirical FDR remains at or below the target~$\alpha$. Figure~\ref{fig:fdr_all} tests this across all candidate--judge--dataset configurations. In every case, the FDR stays at or below~$\alpha$, satisfying the finite-sample guarantee (Eq.~\ref{eq:guarantee_route}).

Three observations are noteworthy. First, all four judges provide valid FDR control regardless of their parameter count, showing that the guarantee does not require a strong judge. Second, FDR tracks close to~$\alpha$ with narrow confidence bands ($\pm$0.01--0.02), which indicates stable calibration across 100 random splits. Third, the gap between FDR and~$\alpha$ is smallest on TriviaQA and PopQA, which are relatively easy tasks where judges are most accurate. This demonstrates that the framework is tightest when the judge has sufficient knowledge to evaluate.

\subsection{Coverage scales with risk tolerance and judge capability}
\label{sec:coverage_results}

Table~\ref{tab:coverage} reports coverage across five risk levels. Coverage scales smoothly with~$\alpha$ and is governed by two factors: dataset difficulty and judge capability. On easy benchmarks (TriviaQA, PopQA), even the smallest judge reaches near-full coverage by $\alpha\!=\!0.25$, while on knowledge-intensive tasks (NQ-Open, HotpotQA) only the larger judges (Qwen-14B, Llama-8B) surpass 80\% at $\alpha\!=\!0.20$ for the Qwen-8B candidate --- Qwen-4B remains below 15\% on these same benchmarks. Evaluating the stronger Llama-70B candidate is consistently harder: coverage drops across all judges compared to Qwen-8B, with the best judge reaching only 53\% on HotpotQA at $\alpha\!=\!0.20$. This gap reflects both the answer diversity of the 70B model and the increased knowledge demands it places on the judge.

\begin{table*}[h]
\centering
\caption{Coverage of the joint two-mode framework (two-threshold routing with retrieval) at $\alpha \in \{.05, .10, .15, .20, .25\}$ (mean over 100 splits) across all datasets.}
\label{tab:coverage}
\resizebox{\textwidth}{!}{%
\begin{tabular}{ll|ccccc|ccccc|ccccc|ccccc}
\toprule
 & & \multicolumn{5}{c|}{TriviaQA} & \multicolumn{5}{c|}{NQ-Open} & \multicolumn{5}{c|}{PopQA} & \multicolumn{5}{c}{HotpotQA} \\
Cand. & Judge & .05 & .10 & .15 & .20 & .25 & .05 & .10 & .15 & .20 & .25 & .05 & .10 & .15 & .20 & .25 & .05 & .10 & .15 & .20 & .25 \\
\midrule
\multirow{4}{*}{\rotatebox{90}{\scriptsize Qwen-8B}}
 & Qwen-4B   & .03 & .39 & .66 & .91 & 1.0 & .00 & .03 & .08 & .14 & .28 & .00 & .01 & .91 & 1.0 & 1.0 & .00 & .04 & .09 & .14 & .30 \\
 & Qwen-8B   & .12 & .63 & .91 & 1.0 & 1.0 & .00 & .02 & .24 & .82 & 1.0 & .00 & .87 & 1.0 & 1.0 & 1.0 & .00 & .04 & .49 & .74 & .96 \\
 & Qwen-14B  & .21 & .78 & .98 & 1.0 & 1.0 & .00 & .05 & .56 & .81 & .99 & .19 & .89 & 1.0 & 1.0 & 1.0 & .02 & .26 & .64 & .86 & 1.0 \\
 & Llama-8B  & .13 & .61 & .86 & 1.0 & 1.0 & .00 & .20 & .61 & .83 & 1.0 & .02 & .88 & 1.0 & 1.0 & 1.0 & .01 & .30 & .65 & .85 & 1.0 \\
\midrule
\multirow{4}{*}{\rotatebox{90}{\scriptsize Llama-70B}}
 & Qwen-4B   & .05 & .37 & .80 & .98 & 1.0 & .00 & .05 & .18 & .32 & .48 & .00 & .00 & .11 & .35 & .92 & .00 & .09 & .20 & .34 & .55 \\
 & Qwen-8B   & .17 & .44 & .83 & .99 & 1.0 & .01 & .10 & .25 & .47 & .73 & .00 & .06 & .11 & .48 & .98 & .00 & .11 & .31 & .47 & .65 \\
 & Qwen-14B  & .12 & .46 & .85 & 1.0 & 1.0 & .00 & .05 & .23 & .51 & .75 & .00 & .04 & .31 & .77 & .99 & .00 & .07 & .31 & .53 & .76 \\
 & Llama-8B  & .04 & .32 & .67 & .86 & 1.0 & .00 & .02 & .13 & .32 & .59 & .00 & .07 & .23 & .77 & .98 & .00 & .02 & .11 & .27 & .51 \\
\bottomrule
\end{tabular}%
}
\end{table*}

\subsection{Adaptive retrieval recovers substantial coverage}
\label{sec:routing_results}
Figure~\ref{fig:routing} breaks down accepted instances by routing mode. It is depicted that retrieval is the dominant acceptance mechanism. Across various configurations, Mode~2 accounts for the majority of non-abstained verdicts. This does not necessarily mean the judge lacks the knowledge to evaluate correctly, but rather that its uncertainty signal in Mode~1 is not confident enough to satisfy the calibrated threshold. Crucially, these are instances where a single-mode framework would abstain entirely. By grounding the judge in retrieved evidence, Mode~2 reduces uncertainty below the second threshold and converts these abstentions into accepted verdicts. This trade-off also depends on the underlying judge model. Models with better-calibrated uncertainty estimates allow a larger fraction of instances to be resolved by Mode~1 directly. For instance, at $\alpha\!=\!0.20$ on TriviaQA, Qwen3-14B resolves 45\% of evaluations parametrically versus 9\% for Qwen3-4B. On harder tasks like HotpotQA and NQ-Open, all judges rely predominantly on retrieved evidence. 

Although retrieval substantially expands coverage, it does not always improve the judge's confidence. For 14--29\% of instances, retrieval increases rather than decreases uncertainty (Appendix~\ref{app:deltau}). The second calibrated threshold identifies these cases and routes them to abstention. Consequently, only retrieval-based verdicts that satisfy the desired risk level are accepted.

\begin{figure*}[htbp]
\centering
\includegraphics[width=\textwidth]{figures/routing_all.pdf}
\caption{Routing distribution across configurations. Within each group of four bars, judges are ordered left to right: Qwen3-4B, Qwen3-8B, Qwen3-14B, Llama-3.1-8B.}
\label{fig:routing}
\end{figure*}

\subsection{Single-mode vs.\ joint two-mode routing}
\label{sec:mode1_vs_joint}

Table~\ref{tab:mode1_vs_joint} isolates the contribution of adaptive retrieval by comparing Mode~1 Only against the full joint framework. On TriviaQA, Mode~1 already achieves high coverage for strong judges (87\% for Qwen3-14B), so retrieval provides a modest but consistent gain. The effect is more pronounced: on NQ-Open, coverage for Qwen3-8B increases from 7\% to 82\%; on HotpotQA, coverage for Llama-8B increases from 40\% to 85\%. Similarly, the weakest judge (Qwen3-4B) benefits substantially on PopQA, with coverage increasing from 4\% to 100\%.


\begin{table*}[!ht]
\centering
\caption{Coverage at $\alpha\!=\!0.20$: Mode~1 Only vs.\ Joint (two-threshold routing with retrieval).}
\label{tab:mode1_vs_joint}
\resizebox{\textwidth}{!}{%
\footnotesize
\begin{tabular}{ll|cc|cc|cc|cc}
\toprule
 & & \multicolumn{2}{c|}{TriviaQA} & \multicolumn{2}{c|}{NQ-Open} & \multicolumn{2}{c|}{PopQA} & \multicolumn{2}{c}{HotpotQA} \\
Cand. & Judge & M1 Only & Joint & M1 Only & Joint & M1 Only & Joint & M1 Only & Joint \\
\midrule
\multirow{4}{*}{\rotatebox{90}{\scriptsize Qwen-8B}}
 & Qwen-4B   & .47 & .91 & .01 & .14 & .04 & 1.0 & .00 & .14 \\
 & Qwen-8B   & .59 & 1.0 & .07 & .82 & .03 & 1.0 & .05 & .74 \\
 & Qwen-14B  & .87 & 1.0 & .14 & .81 & .15 & 1.0 & .35 & .86 \\
 & Llama-8B  & .77 & 1.0 & .22 & .83 & .21 & 1.0 & .40 & .85 \\
\midrule
\multirow{4}{*}{\rotatebox{90}{\scriptsize Llama-70B}}
 & Qwen-4B   & .67 & .98 & .07 & .32 & .05 & .35 & .09 & .34 \\
 & Qwen-8B   & .82 & .99 & .22 & .47 & .08 & .48 & .27 & .47 \\
 & Qwen-14B  & .84 & 1.0 & .12 & .51 & .15 & .77 & .24 & .53 \\
 & Llama-8B  & .48 & .86 & .02 & .32 & .12 & .77 & .02 & .27 \\
\bottomrule
\end{tabular}%
}
\end{table*}
\subsection{Robustness to retrieval non-stationarity}
\label{sec:drift}
Web search results change over time, which \S\ref{sec:guarantee} identified as the main practical threat to the determinism the guarantee assumes for Mode~2. We measure the effect directly on the main configuration (Qwen3-8B candidate, Qwen3-8B judge, LLM evaluator, $k\!=\!3$). We re-query the same 2{,}000 questions on TriviaQA and HotpotQA roughly three months after the original snapshot and recompute the Mode~2 verdicts on the fresh evidence. The retrieved pages turn over substantially. The mean per-item Jaccard overlap between the original ($\tau_0$) and re-queried ($\tau_1$) ranked top-$3$ URL lists is $0.30$ for TriviaQA and $0.24$ for HotpotQA. When ignoring the retrieval rank and considering only URL membership, the corresponding position-blind top-$3$ overlaps are $0.47$ and $0.40$, respectively.

To evaluate whether the calibrated thresholds remain valid after retrieval drift, we calibrate $(\hat{t}_1, \hat{t}_2)$ on the original data and apply them to test instances evaluated with the updated Mode~2 verdicts (while keeping Mode~1 unchanged). We repeat this evaluation over 100 random 50/50 splits. Table~\ref{tab:drift} summarizes the results. Across all settings, the empirical FDR remains below the target level~$\alpha$. Similarly, coverage remains largely stable after redeployment. This indicates that drift changes which pages are returned more than the overall quality of the evidence. Therefore, the calibrated thresholds remain effective after redeployment.

\begin{table*}[!ht]
\centering
\footnotesize
\caption{FDR and coverage under retrieval drift from the original
($\tau_0$) to the re-queried snapshot ($\tau_1$, $\sim$3 months later).
$\Delta=\tau_1-\tau_0$.}
\label{tab:drift}
\begin{tabular*}{\textwidth}{@{\extracolsep{\fill}}c|ccc|ccc|ccc|ccc@{}}
\toprule
 & \multicolumn{6}{c|}{TriviaQA} & \multicolumn{6}{c}{HotpotQA} \\
 & \multicolumn{3}{c|}{FDR} & \multicolumn{3}{c|}{Coverage} & \multicolumn{3}{c|}{FDR} & \multicolumn{3}{c}{Coverage} \\
$\alpha$ & $\tau_0$ & $\tau_1$ & $\Delta$ & $\tau_0$ & $\tau_1$ & $\Delta$ & $\tau_0$ & $\tau_1$ & $\Delta$ & $\tau_0$ & $\tau_1$ & $\Delta$ \\
\midrule
.05 & .03 & .04 & .00 & .12 & .17 & $+$.05 & .00 & .00 & .00 & .00 & .00 & .00 \\
.10 & .08 & .09 & .00 & .63 & .68 & $+$.06 & .05 & .03 & $-$.01 & .04 & .02 & $-$.02 \\
.15 & .14 & .13 & .00 & .91 & .94 & $+$.04 & .13 & .13 & .00 & .49 & .49 & .00 \\
.20 & .16 & .15 & $-$.01 & 1.0 & 1.0 & .00 & .18 & .18 & .00 & .74 & .72 & $-$.02 \\
.25 & .16 & .15 & $-$.01 & 1.0 & 1.0 & .00 & .23 & .23 & .00 & .96 & .96 & .00 \\
\bottomrule
\end{tabular*}
\end{table*}

% \subsection{Robustness to retrieval non-stationarity}
% \label{sec:drift}
% Web search results change over time, which \S\ref{sec:guarantee} identified as the main practical threat to the determinism the guarantee assumes for Mode~2. We measure the effect directly on the main configuration (Qwen3-8B candidate, Qwen3-8B judge, LLM evaluator, $k\!=\!3$). We re-query the same 2{,}000 questions on TriviaQA and HotpotQA roughly three months after the original snapshot and recompute the Mode~2 verdicts on the fresh evidence. The retrieved pages turn over substantially. The mean per-item Jaccard overlap between the old and new ranked top-$3$ URL lists is $0.30$ for TriviaQA and $0.24$ for HotpotQA. When ignoring the retrieval rank and considering only URL membership, the corresponding position-blind top-$3$ overlaps are $0.47$ and $0.40$, respectively.

% To evaluate whether the calibrated thresholds remain valid after retrieval drift, we calibrate $(\hat{t}_1, \hat{t}_2)$ on the original data and apply them to test instances evaluated with the updated Mode~2 verdicts (while keeping Mode~1 unchanged). We repeat this evaluation over 100 random 50/50 splits. Table~\ref{tab:drift} summarizes the results. Across all settings, the empirical FDR remains below the target level~$\alpha$, with no observed violations. The absolute change in FDR is at most $0.012$. Similarly, coverage remains largely stable after redeployment. It increases on TriviaQA (up to $+0.057$ at $\alpha=0.10$) and decreases slightly on HotpotQA at larger~$\alpha$ values (up to $-0.021$). These results suggest that retrieval drift primarily changes the specific pages returned rather than the overall quality of the retrieved evidence. As a result, the calibrated thresholds remain effective under the updated retrieval conditions. Nevertheless, recalibration may be necessary when the underlying evidence distribution changes substantially.

% \begin{table}[t]
% \centering
% \small
% \caption{FDR and Coverage under retrieval drift. Thresholds are calibrated on the original snapshot and applied to test items re-scored with evidence re-queried $\sim$3 months later.}
% \label{tab:drift}
% \begin{tabular}{ll|ccccc}
% \toprule
%  & & \multicolumn{5}{c}{$\alpha$} \\
% Dataset & Metric & .05 & .10 & .15 & .20 & .25 \\
% \midrule
% \multirow{4}{*}{TriviaQA}
%  & FDR (old)   & .03 & .08 & .14 & .16 & .16 \\
%  & FDR (new)   & .04 & .09 & .13 & .15 & .15 \\
%  & Cov.\ (old) & .12 & .63 & .91 & 1.0 & 1.0 \\
%  & Cov.\ (new) & .17 & .68 & .94 & 1.0 & 1.0 \\
% \midrule
% \multirow{4}{*}{HotpotQA}
%  & FDR (old)   & .00 & .05 & .13 & .18 & .23 \\
%  & FDR (new)   & .00 & .03 & .13 & .18 & .23 \\
%  & Cov.\ (old) & .00 & .04 & .49 & .74 & .96 \\
%  & Cov.\ (new) & .00 & .02 & .49 & .72 & .96 \\
% \bottomrule
% \end{tabular}
% \end{table}

